\section{第五章：一般理论与近似计算}
\basis{第5.1--5.6节，教材第114--159页。}
\problem{5.1.1}{121}{在矩形 $|t-t_0|\le a,|x-x_0|\le b$ 上，$f$ 连续且关于 $x$ 的 Lipschitz 常数为 $L$。设 $M=\max|f|$，$h=\min(a,b/M)$，$\varphi$ 是初值解，$\varphi_n$ 是从 $\varphi_0=x_0$ 开始的第 $n$ 次 Picard 迭代。证明误差估计。}
\begin{proof}
记误差 $e_n(t)=|\varphi_n(t)-\varphi(t)|$。从积分方程和 $\varphi_0=x_0$ 出发，先估计第0步：
\[
e_0(t)=\left|\int_{t_0}^tf(s,\varphi(s))\dd s\right|
\le M|t-t_0|.
\]
第1步则满足 $\displaystyle e_1(t)\le L\int_{t_0}^te_0(s)\dd s\le ML(t-t_0)^2/2$，此处先写 $t\ge t_0$。

若 $e_{n-1}(s)\le ML^{n-1}(s-t_0)^n/n!$，由两份积分式相减并使用 Lipschitz 条件，
\[
\begin{aligned}
e_n(t)&\le\int_{t_0}^t|f(s,\varphi_{n-1}(s))-f(s,\varphi(s))|\dd s\\
&\le L\int_{t_0}^te_{n-1}(s)\dd s\\
&\le\dfrac{ML^n}{n!}\int_{t_0}^t(s-t_0)^n\dd s
=\dfrac{ML^n(t-t_0)^{n+1}}{(n+1)!}.
\end{aligned}
\]
这证明了归纳步。若 $t<t_0$，把取绝对值后的积分改为 $\displaystyle \int_t^{t_0}$，以 $|s-t_0|$ 替代 $s-t_0$，计算结果为同一绝对值公式。最后在 $|t-t_0|\le h$ 上取上界。
\ans{|\varphi_n(t)-\varphi(t)|\le\dfrac{ML^nh^{n+1}}{(n+1)!},\quad |t-t_0|\le h.}
$L=0$ 时一次迭代就已精确，另行按此解释，不需使用 $0^0$。
\end{proof}
\problem{5.1.2}{121}{证明 $x'=Ax,x(0)=x_0$ 的 Picard 迭代收敛到 $\exp(At)x_0$，$A$ 为常矩阵。}
Picard 迭代从常函数 $\varphi_0(t)=x_0$ 开始，递推为 $\displaystyle \varphi_{n+1}(t)=x_0+\int_0^tA\varphi_n(s)\dd s$。先算前两步：
\[
\begin{aligned}
\varphi_1&=x_0+\int_0^tAx_0\dd s=(I+At)x_0,\\
\varphi_2&=x_0+\int_0^t[A+A^2s]x_0\dd s
=\left(I+At+\dfrac{A^2t^2}{2}\right)x_0.
\end{aligned}
\]
若第 $n$ 步为 $\sum_{k=0}^nA^kt^kx_0/k!$，则代回积分式：
\[
\begin{aligned}
\varphi_{n+1}
&=x_0+\sum_{k=0}^n\dfrac{A^{k+1}x_0}{k!}\int_0^t s^k\dd s\\
&=x_0+\sum_{k=0}^n\dfrac{A^{k+1}t^{k+1}}{(k+1)!}x_0
=\sum_{j=0}^{n+1}\dfrac{A^jt^j}{j!}x_0.
\end{aligned}
\]
故通式由归纳法成立。在 $|t|\le T$ 上，第 $k$ 项范数不超过 $\|x_0\|\|A\|^kT^k/k!$，控制级数为收敛的指数级数，因此可一致取极限，得到矩阵指数定义。
\ans{x(t)=\e^{At}x_0.}
级数逐项求导得到 $x'=Ax$，在0处为 $x_0$，由唯一性即为初值解。
\problem{5.1.3}{121}{求 $x'=x^2,x(0)=1$ 的第三次 Picard 近似解。}
按 $\varphi_0=1$、$\displaystyle \varphi_{n+1}=1+\int_0^t\varphi_n(s)^2\dd s$ 逐步计算。第一次：
\[
\varphi_1=1+\int_0^t1\dd s=1+t.
\]
第二次先平方，再积分：
\[
\varphi_2=1+\int_0^t(1+2s+s^2)\dd s
=1+t+t^2+\dfrac{t^3}{3}.
\]
第三次需要把四项多项式平方完整展开：
\[
\begin{aligned}
\left(1+s+s^2+\dfrac{s^3}{3}\right)^2
&=1+2s+3s^2+\dfrac83s^3+\dfrac53s^4+\dfrac23s^5+\dfrac19s^6.
\end{aligned}
\]
因此 $\varphi_3$ 是1加上该多项式从0到 $t$ 的积分；各系数依次除以新的幂次，得到 $1,1,1,2/3,1/3,1/9,1/63$，如下所列。不要把精确解的三阶 Taylor 多项式当作“第三次迭代”，两者不是同一对象。
\ans{\varphi_3=1+t+t^2+t^3+\dfrac{2t^4}{3}+\dfrac{t^5}{3}+\dfrac{t^6}{9}+\dfrac{t^7}{63}.}
\checknote{$\varphi_3(0)=1$ 且 $\varphi_3'=\varphi_2^2$；精确解 $1/(1-t)$ 的前四项与之相同，最大存在区间为 $(-\infty,1)$。}
\problem{5.1.4}{121}{对矩形 $|t|\le1,|x|\le1$ 内的方程 $x'=x^2-t$，求过 $(0,0)$ 的解由 Picard 定理保证的存在区间、第三次近似解及该区间上的误差估计。}
先估计矩形内的常数。$|x^2-t|\le1+1=2$，且在 $(t,x)=(-1,\pm1)$ 取到2，故 $M=2$；由 $|f_x|=|2x|\le2$，可取 $L=2$。矩形半宽和半高均为1，所以
\[
h=\min\left(1,\dfrac12\right)=\dfrac12.
\]
从 $\varphi_0=0$ 开始，递推为 $\displaystyle \varphi_{n+1}=\int_0^t[\varphi_n(s)^2-s]\dd s$。逐步得到
\[
\begin{aligned}
\varphi_1&=\int_0^t(-s)\dd s=-\dfrac{t^2}{2},\\
\varphi_2&=\int_0^t\left(\dfrac{s^4}{4}-s\right)\dd s
=-\dfrac{t^2}{2}+\dfrac{t^5}{20}.
\end{aligned}
\]
第三次的被积函数为
\[
\left(-\dfrac{s^2}{2}+\dfrac{s^5}{20}\right)^2-s
=\dfrac{s^4}{4}-\dfrac{s^7}{20}+\dfrac{s^{10}}{400}-s.
\]
逐项积分得到下面的 $\varphi_3$。误差公式取 $n=3$，将 $M=L=2,h=1/2$ 逐一代入：$ML^3h^4/4!=2\cdot8\cdot(1/16)/24=1/24$。
\ans{I=[-1/2,1/2],\qquad
\varphi_3=-\dfrac{t^2}{2}+\dfrac{t^5}{20}-\dfrac{t^8}{160}+\dfrac{t^{11}}{4400},\\
|\varphi_3-\varphi|\le\dfrac{2\cdot2^3(1/2)^4}{4!}=\dfrac1{24}.}
这里 $I$ 是该矩形给出的保证区间，不代表真实最大存在区间。
\problem{5.1.5}{121}{对同一矩形内的 $x'=x^2+t,x(0)=0$，求保证存在区间及在此区间误差不超过0.05的近似解。}
矩形上 $|x^2+t|\le2$ 且取到2，故 $M=2$；$|f_x|\le2$，取 $L=2$。由 $h=\min(1,1/2)=1/2$ 得保证区间。为选迭代次数，先列通用误差界：
\[
E_n=\dfrac{ML^nh^{n+1}}{(n+1)!}
=\dfrac{2\cdot2^n(1/2)^{n+1}}{(n+1)!}=\dfrac1{(n+1)!}.
\]
$n=2$ 时界为 $1/6>0.05$；$n=3$ 时为 $1/24<0.05$，所以计算到第三次足够。

由 $\displaystyle \varphi_{n+1}=\int_0^t[\varphi_n(s)^2+s]\dd s$、$\varphi_0=0$，
\[
\varphi_1=\dfrac{t^2}{2},\qquad
\varphi_2=\int_0^t\left(\dfrac{s^4}{4}+s\right)\dd s
=\dfrac{t^2}{2}+\dfrac{t^5}{20}.
\]
第三次先展开平方：$\varphi_2(s)^2=s^4/4+s^7/20+s^{10}/400$，再加 $s$ 逐项积分，得到下式。
\ans{I=[-1/2,1/2],\qquad
\varphi_3=\dfrac{t^2}{2}+\dfrac{t^5}{20}+\dfrac{t^8}{160}+\dfrac{t^{11}}{4400}.}
第二次近似的通用误差界为 $1/6$，不足以用此估计保证题定精度；这不表示实际误差一定大于0.05。
\problem{5.1.6}{121--122}{求 $x'=x+t+1,x(0)=0$ 的 Picard 迭代序列，并由其极限求解。}
由积分式 $\displaystyle x(t)=\int_0^t[x(s)+s+1]\dd s$，取 $\varphi_0=0$，逐次代入：
\[
\begin{aligned}
\varphi_1&=\int_0^t(s+1)\dd s=t+\dfrac{t^2}{2},\\
\varphi_2&=\int_0^t\left(1+2s+\dfrac{s^2}{2}\right)\dd s
=t+t^2+\dfrac{t^3}{6},\\
\varphi_3&=\int_0^t\left(1+2s+s^2+\dfrac{s^3}{6}\right)\dd s
=t+t^2+\dfrac{t^3}{3}+\dfrac{t^4}{24}.
\end{aligned}
\]
这提示 $\varphi_n=2\sum_{k=0}^nt^k/k!-t-2+t^{n+1}/(n+1)!$。为验证通式，将它代入下一次的被积函数，$-s$ 与外加的 $s$ 抵消、$-2$ 与外加1合并为 $-1$。积分后成为
\[
2\sum_{k=0}^n\dfrac{t^{k+1}}{(k+1)!}-t+\dfrac{t^{n+2}}{(n+2)!},
\]
这正是通式的 $n+1$ 形式。有限区间上指数级数一致收敛，末尾单项趋零，故可取极限。
\ans{\varphi_n=2\sum_{k=0}^n\dfrac{t^k}{k!}-t-2+\dfrac{t^{n+1}}{(n+1)!},\quad
x=\lim_{n\to\infty}\varphi_n=2\e^t-t-2.}
\checknote{极限满足 $x'=2\e^t-1=x+t+1$ 且 $x(0)=0$；有限区间上一致收敛。}
\problem{5.1.7}{122}{求 $x'=P(t)x+Q(t),x(t_0)=x_0$ 的 Picard 序列及其极限，其中 $P,Q$ 连续。}
设 $\displaystyle A(t)=\int_{t_0}^tP(s)\dd s$，并取 $\varphi_0=x_0$。第一步为
\[
\varphi_1=x_0+x_0A(t)+\int_{t_0}^tQ(s)\dd s.
\]
第二步代入时，$P(s)A(s)$ 的积分为 $A(t)^2/2$；对 $Q$ 产生的重积分交换次序，得
\[
\varphi_2=x_0\left[1+A(t)+\dfrac{A(t)^2}{2}\right]
+\int_{t_0}^tQ(s)[1+A(t)-A(s)]\dd s.
\]
这两类项分别由初值和外力 $Q$ 产生。归纳式为
\[
\varphi_n=x_0\sum_{j=0}^n\dfrac{A(t)^j}{j!}
+\int_{t_0}^tQ(s)\sum_{j=0}^{n-1}\dfrac{[A(t)-A(s)]^j}{j!}\dd s.
\]
归纳一步对每个 $s$ 使用 $\displaystyle \int_s^tP(u)[A(u)-A(s)]^j\dd u=[A(t)-A(s)]^{j+1}/(j+1)$。在任意有限闭区间上相关连续函数有界，指数级数一致收敛，故可以在积分号内取极限：
\[
\varphi_n\longrightarrow x_0\e^{A(t)}+
\int_{t_0}^tQ(s)\e^{A(t)-A(s)}\dd s.
\]
提出只依赖 $t$ 的 $\e^{A(t)}$，得到标准的线性方程求解公式。
\ans{x(t)=\e^{A(t)}\left[x_0+\int_{t_0}^t\e^{-A(s)}Q(s)\dd s\right].}
对该式求导并代入 $t_0$，也独立恢复线性方程及初值。
\problem{5.1.8}{122}{用逐步逼近证明：$|\lambda|$ 充分小时，积分方程 $\varphi(t)=f(t)+\lambda\displaystyle\int_a^bK(t,s)\varphi(s)\dd s$ 在 $[a,b]$ 上有唯一连续解，其中 $f,K$ 连续。}
\begin{proof}
设 $M=\max|K|$，取 $q=|\lambda|M(b-a)<1$。令 $(Tu)(t)=\displaystyle\int_a^bK(t,s)u(s)\dd s$，则 $\|Tu\|_\infty\le M(b-a)\|u\|_\infty$。
从 $\varphi_0=f$ 出发，令 $\varphi_{n+1}=f+\lambda T\varphi_n$。由归纳
\[
 \varphi_n=\sum_{j=0}^n\lambda^jT^jf,\qquad
 \|\varphi_{n+1}-\varphi_n\|_\infty\le q^{n+1}\|f\|_\infty.
\]
几何级数收敛，故序列一致收敛到连续函数 $\varphi$。$T$ 的上述范数界保证可把极限移入积分，从而 $\varphi=f+\lambda T\varphi$。
若另有解 $\psi$，则 $\|\varphi-\psi\|\le q\|\varphi-\psi\|$。因为 $q<1$，必有 $\varphi=\psi$。若 $M=0$，方程直接成为 $\varphi=f$，任意 $\lambda$ 都可。
收敛的定量理由是对 $m>n$ 有 $\|\varphi_m-\varphi_n\|_\infty\le\|f\|_\infty\sum_{j=n+1}^mq^j\le\|f\|_\infty q^{n+1}/(1-q)$，它与 $m$ 无关且趋零。因此一致 Cauchy 性并非只由相邻差趋零得到，而是由相邻差的可求和几何上界得到。
\end{proof}
\problem{5.2.1}{128}{对 $x'=x,x(0)=1$，把 $[-1,1]$ 分为8等份，构造 Euler 折线，与精确曲线比较。计算各等分点近似值及误差，精确到小数点后第五位。}
八等分给出节点 $t_k=k/4$，$k=-4,\ldots,4$，初值位于中间的 $t_0=0$，不是左端点。精确解由分离变量为 $x(t)=\e^t$。

正向步长 $h=1/4$，Euler 公式为 $X_{k+1}=X_k+hX_k=(5/4)X_k$。从 $X_0=1$ 出发，依次为
\[
X_1=1.25,\quad X_2=1.5625,\quad X_3=1.953125,\quad X_4=2.44140625.
\]
反向仍从0处初值出发，但取步长 $-1/4$：
\[
X_{-(k+1)}=X_{-k}-\dfrac14X_{-k}=\dfrac34X_{-k}.
\]
因此 $X_{-1}=0.75,X_{-2}=0.5625,X_{-3}=0.421875,X_{-4}=0.31640625$。

对每一步，以推进的起点 $(t_j,X_j)$ 为基点画线段 $\varphi(t)=X_j+(t-t_j)X_j$。例如正向首段为 $1+t$（$0\le t\le1/4$），反向首段也是 $1+t$（$-1/4\le t\le0$）；后续负向区间应取其右端点作为推进起点。节点误差逐项计算 $|\e^{t_k}-X_k|$，最后才四舍五入到五位小数。
\ans{X_k=(5/4)^k\ (k=0,1,\ldots,4),\qquad
 X_{-k}=(3/4)^k\ (k=0,1,\ldots,4).}
在相邻节点之间用 $\varphi(t)=X_j+(t-t_j)X_j$（沿相应积分方向选择起点 $t_j$）作直线段。精确解为 $\e^t$。
\begin{center}\begin{tabular}{rrrr}\toprule $t$&$\e^t$&Euler近似&绝对误差\\\midrule
-1.00&0.36788&0.31641&0.05147\\
-0.75&0.47237&0.42188&0.05049\\
-0.50&0.60653&0.56250&0.04403\\
-0.25&0.77880&0.75000&0.02880\\
0.00&1.00000&1.00000&0.00000\\
0.25&1.28403&1.25000&0.03403\\
0.50&1.64872&1.56250&0.08622\\
0.75&2.11700&1.95312&0.16388\\
1.00&2.71828&2.44141&0.27688\\
\bottomrule\end{tabular}\end{center}
\begin{center}\begin{tikzpicture}[x=3.5cm,y=1.25cm]
\draw[->] (-1.1,0)--(1.15,0) node[right]{$t$};
\draw[->] (0,0)--(0,3) node[above]{$x$};
\draw[blue,thick] plot coordinates {(-1.00000,0.36788) (-0.98658,0.37285) (-0.97315,0.37789) (-0.95973,0.38300) (-0.94631,0.38817) (-0.93289,0.39342) (-0.91946,0.39873) (-0.90604,0.40412) (-0.89262,0.40958) (-0.87919,0.41512) (-0.86577,0.42073) (-0.85235,0.42641) (-0.83893,0.43217) (-0.82550,0.43801) (-0.81208,0.44393) (-0.79866,0.44993) (-0.78523,0.45601) (-0.77181,0.46217) (-0.75839,0.46842) (-0.74497,0.47475) (-0.73154,0.48117) (-0.71812,0.48767) (-0.70470,0.49426) (-0.69128,0.50094) (-0.67785,0.50771) (-0.66443,0.51457) (-0.65101,0.52152) (-0.63758,0.52857) (-0.62416,0.53571) (-0.61074,0.54295) (-0.59732,0.55029) (-0.58389,0.55772) (-0.57047,0.56526) (-0.55705,0.57290) (-0.54362,0.58064) (-0.53020,0.58849) (-0.51678,0.59644) (-0.50336,0.60450) (-0.48993,0.61267) (-0.47651,0.62095) (-0.46309,0.62934) (-0.44966,0.63784) (-0.43624,0.64646) (-0.42282,0.65520) (-0.40940,0.66405) (-0.39597,0.67302) (-0.38255,0.68212) (-0.36913,0.69134) (-0.35570,0.70068) (-0.34228,0.71015) (-0.32886,0.71974) (-0.31544,0.72947) (-0.30201,0.73933) (-0.28859,0.74932) (-0.27517,0.75944) (-0.26174,0.76971) (-0.24832,0.78011) (-0.23490,0.79065) (-0.22148,0.80133) (-0.20805,0.81216) (-0.19463,0.82314) (-0.18121,0.83426) (-0.16779,0.84554) (-0.15436,0.85696) (-0.14094,0.86854) (-0.12752,0.88028) (-0.11409,0.89217) (-0.10067,0.90423) (-0.08725,0.91645) (-0.07383,0.92883) (-0.06040,0.94139) (-0.04698,0.95411) (-0.03356,0.96700) (-0.02013,0.98007) (-0.00671,0.99331) (0.00671,1.00673) (0.02013,1.02034) (0.03356,1.03413) (0.04698,1.04810) (0.06040,1.06226) (0.07383,1.07662) (0.08725,1.09117) (0.10067,1.10591) (0.11409,1.12086) (0.12752,1.13600) (0.14094,1.15136) (0.15436,1.16691) (0.16779,1.18268) (0.18121,1.19866) (0.19463,1.21486) (0.20805,1.23128) (0.22148,1.24792) (0.23490,1.26478) (0.24832,1.28187) (0.26174,1.29920) (0.27517,1.31675) (0.28859,1.33455) (0.30201,1.35258) (0.31544,1.37086) (0.32886,1.38938) (0.34228,1.40816) (0.35570,1.42719) (0.36913,1.44647) (0.38255,1.46602) (0.39597,1.48583) (0.40940,1.50591) (0.42282,1.52626) (0.43624,1.54688) (0.44966,1.56779) (0.46309,1.58897) (0.47651,1.61044) (0.48993,1.63221) (0.50336,1.65426) (0.51678,1.67662) (0.53020,1.69927) (0.54362,1.72224) (0.55705,1.74551) (0.57047,1.76910) (0.58389,1.79300) (0.59732,1.81723) (0.61074,1.84179) (0.62416,1.86668) (0.63758,1.89190) (0.65101,1.91747) (0.66443,1.94338) (0.67785,1.96964) (0.69128,1.99626) (0.70470,2.02324) (0.71812,2.05058) (0.73154,2.07829) (0.74497,2.10637) (0.75839,2.13483) (0.77181,2.16368) (0.78523,2.19292) (0.79866,2.22256) (0.81208,2.25259) (0.82550,2.28303) (0.83893,2.31388) (0.85235,2.34515) (0.86577,2.37684) (0.87919,2.40896) (0.89262,2.44151) (0.90604,2.47450) (0.91946,2.50794) (0.93289,2.54183) (0.94631,2.57618) (0.95973,2.61100) (0.97315,2.64628) (0.98658,2.68204) (1.00000,2.71828)};
\draw[red,thick] plot coordinates {(-1.00000,0.31641) (-0.75000,0.42188) (-0.50000,0.56250) (-0.25000,0.75000) (0.00000,1.00000) (0.25000,1.25000) (0.50000,1.56250) (0.75000,1.95312) (1.00000,2.44141)};
\fill[red] (-1.00000,0.31641) circle (1.1pt);
\fill[red] (-0.75000,0.42188) circle (1.1pt);
\fill[red] (-0.50000,0.56250) circle (1.1pt);
\fill[red] (-0.25000,0.75000) circle (1.1pt);
\fill[red] (0.00000,1.00000) circle (1.1pt);
\fill[red] (0.25000,1.25000) circle (1.1pt);
\fill[red] (0.50000,1.56250) circle (1.1pt);
\fill[red] (0.75000,1.95312) circle (1.1pt);
\fill[red] (1.00000,2.44141) circle (1.1pt);
\node[blue] at (.65,2.8){$\e^t$};
\node[red] at (-.65,1.1){Euler折线};
\end{tikzpicture}\end{center}
\checknote{所有节点均从给定 $X_0=1$ 递推，未把 $x(-1)=\e^{-1}$ 当作另一个精确起点。误差表用未提前舍入的数据计算。}
\problem{5.3.1}{135}{设 $G=\{(t,x):T_0<t<T_1,|x|<\infty\}$，$f$ 在 $G$ 连续。若 $x'=f(t,x)$ 的一个饱和解有界，证明其最大存在区间为 $(T_0,T_1)$。}
\begin{proof}
设最大区间为 $(\alpha,\beta)$。若 $\beta<T_1$ 且为有限内点，取 $c\in(\alpha,\beta)$。由解有界，其图像在 $[c,\beta)$ 上落在 $[c,\beta]\times[-K,K]$ 这个 $G$ 内的紧矩形中，$|f|$ 有界，故 $x'$ 有界。于是 $t\uparrow\beta$ 时 $x(t)$ 为 Cauchy，存在有限极限 $\xi$。

点 $(\beta,\xi)\in G$，由 Peano 定理可从该点继续解；与原解接合后，积分式保证接合函数可微且仍满足原方程，矛盾。因此右端只能为 $T_1$。对左端完全相同地反向积分，得 $\alpha=T_0$。无限端点按同样的有限内点排除法处理。
端点极限存在可直接由导数界检验：若紧矩形上 $|f|\le M$，则 $|x(t)-x(s)|\le M|t-s|$。当 $s,t$ 同时趋于有限右端 $\beta$ 时，这个差趋零，故 $x(t)$ 的有限极限存在。把极限作为接合点函数值，左右导数均由 $f$ 的连续性趋于 $f(\beta,\xi)$，接合处也满足原方程。
\end{proof}
\problem{5.3.4}{135}{证明 $x'=x^2+t^2$ 的任一解的最大存在区间两端均有限。}
\begin{proof}
若某解能向右存在至无穷，则由 $x'\ge t^2$，
\[
 x(t)\ge x(t_0)+\dfrac{t^3-t_0^3}{3},
\]
所以某个 $T$ 以后 $x>0$。此后 $(1/x)'=-x'/x^2\le-1$，故 $1/x(t)\le1/x(T)-(t-T)$，在 $t>T+1/x(T)$ 时右端为负，与左端正矛盾。因此右端有限。

若能向左存在至负无穷，同一积分不等式说明 $x(t)\to-\infty$。令 $u(s)=-x(-s)$，则 $u'=u^2+s^2$，而 $s\to+\infty$ 时 $u>0$，由前一论证也不可能全局向右存在。因此原解左端也有限。
右向矛盾给出一个具体的存在时间上界：一旦选到 $x(T)>0$，后续寿命至多为 $1/x(T)$。若在这个时长之前解已失去定义，右端自然更早有限；若仍存在到更晚时刻，则倒数不等式产生矛盾。两种情形均排除了右端为无穷。
\end{proof}
\problem{5.3.5(1)}{135}{求 $x'=1+\ln t,x(1)=0$ 的最大存在区间及两端性质。}
原方程的 $\ln t$ 要求 $t>0$。先分部积分：
\[
\int\ln t\dd t=t\ln t-\int1\dd t=t\ln t-t.
\]
所以 $\displaystyle x=\int(1+\ln t)\dd t=t\ln t+C$。由 $x(1)=0$ 得 $C=0$。

所得表达式在所有 $t>0$ 都可微，且右端也在该区间连续，因此最大区间是 $(0,\infty)$。左端极限用 $t\ln t=\ln t/(1/t)$ 求出为0；右端随 $t\to\infty$ 趋于正无穷。尽管左端有有限函数极限，$\ln t$ 仍不能定义于0，故不能据此延拓原方程的解。
\ans{x=t\ln t,\quad (\alpha,\beta)=(0,\infty),\quad
\lim_{t\downarrow0}x(t)=0,\quad\lim_{t\to\infty}x(t)=+\infty.}
$t=0$ 虽有有限函数极限，但原方程右端未定义，不能作为该方程的延拓点。
\problem{5.3.5(2)}{135}{求 $x'=(x^2-1)/2,x(\ln2)=-3$ 的最大存在区间及两端性质。}
初值 $-3$ 不等于平衡点 $\pm1$。分离变量并作部分分式：
\[
\begin{aligned}
\dfrac{2\dd x}{x^2-1}&=\dd t,\\
\ln\left|\dfrac{x-1}{x+1}\right|&=t+C_0.
\end{aligned}
\]
初值点处 $(x-1)/(x+1)=(-4)/(-2)=2>0$，故相应分支写为 $(x-1)/(x+1)=K\e^t$。由 $t=\ln2$ 得 $2=2K$，所以 $K=1$。交叉相乘并移项：
\[
x-1=\e^t(x+1),\qquad (1-\e^t)x=1+\e^t.
\]
分母仅在0处为零，包含初值时刻 $\ln2>0$ 的最大连通区间是 $(0,\infty)$。$t\downarrow0$ 时分子趋2、分母从负侧趋0，故解趋负无穷；$t\to\infty$ 时分子分母同除以 $\e^t$，极限为 $1/(-1)=-1$。
\ans{x=\dfrac{1+\e^t}{1-\e^t},\quad(\alpha,\beta)=(0,\infty),\quad
\lim_{t\downarrow0}x(t)=-\infty,\quad\lim_{t\to\infty}x(t)=-1.}
分母在0处为零；负半轴上同一公式的另一分支不是包含初值点的解。
\problem{5.4.1}{142}{证明 Gronwall 不等式：$x,f$ 在 $[t_0,t_1]$ 连续，$f\ge0$，若 $x(t)\le g+\displaystyle\int_{t_0}^tf(s)x(s)\dd s$，$g$ 为实常数，则 $x(t)\le g\exp(\displaystyle\int_{t_0}^tf(s)\dd s)$。}
\begin{proof}
设 $H(t)=g+\displaystyle\int_{t_0}^tf(s)x(s)\dd s$，则 $x\le H$，并因 $f\ge0$，有 $H'=fx\le fH$。乘积分因子，
\[
 \dfrac\dd{\dd t}\left[H(t)\e^{-\int_{t_0}^tf(s)\dd s}\right]\le0.
\]
从 $t_0$ 到 $t$ 积分，得 $\displaystyle H(t)\le g\e^{\int_{t_0}^tf}$，再用 $x\le H$ 即得结论。此法无需假设 $x$ 或 $g$ 非负。
积分因子的求导为
\[
(H\e^{-A})'=\e^{-A}(H'-A'H)=\e^{-A}(H'-fH)\le0,
\quad A(t)=\int_{t_0}^tf(s)\dd s.
\]
因 $A(t_0)=0$、$H(t_0)=g$，从单调性得 $H(t)\e^{-A(t)}\le g$。再乘恒正的 $\e^{A(t)}$，不论 $g$ 正负，不等号方向都保持不变。
\end{proof}
\problem{5.4.2}{142--143}{设 $f$ 关于 $x$ 的 Lipschitz 常数为 $L$，$\varphi_i$ 为方程 $x'=f(t,x)$ 的 $\varepsilon_i$ 近似解，且 $|\varphi_1(t_0)-\varphi_2(t_0)|\le\delta$。估计两近似解之差，$t\ge t_0$。}
近似解的定义给出残差 $r_i(t)=\varphi_i'(t)-f(t,\varphi_i(t))$，满足 $|r_i|\le\varepsilon_i$。令 $d=\varphi_1-\varphi_2$，由积分式
\[
d(t)=d(t_0)+\int_{t_0}^t[f(s,\varphi_1)-f(s,\varphi_2)+r_1-r_2]\dd s.
\]
取绝对值，使用初始误差界、Lipschitz 条件及 $|r_1-r_2|\le\varepsilon_1+\varepsilon_2=\varepsilon$，得到
\[
u(t):=|d(t)|\le\delta+L\int_{t_0}^tu(s)\dd s+\varepsilon(t-t_0).
\]
把右端记为 $H(t)$，则 $u\le H$、$H(t_0)=\delta$，且 $H'=Lu+\varepsilon\le LH+\varepsilon$。当 $L>0$ 时，乘 $\e^{-L(t-t_0)}$ 并积分：
\[
\begin{aligned}
H(t)\e^{-L(t-t_0)}-\delta
&\le\varepsilon\int_{t_0}^t\e^{-L(s-t_0)}\dd s\\
&=\dfrac\varepsilon L[1-\e^{-L(t-t_0)}].
\end{aligned}
\]
再乘回正指数并用 $u\le H$，得到所需误差界。$L=0$ 时直接由 $H'\le\varepsilon$ 积分，无需作除以 $L$ 的运算。
\ans{u(t)\le\delta\e^{L(t-t_0)}+\dfrac\varepsilon L[\e^{L(t-t_0)}-1]\quad(L>0),\\
u(t)\le\delta+\varepsilon(t-t_0)\quad(L=0).}
$L\to0$ 的极限正是第二式；两近似解为精确解时 $\varepsilon=0$，恢复初值敏感性估计。
\problem{5.4.3}{143}{不用推广的 Gronwall 不等式，证明：若 $f\ge0,K\ge0$ 且 $x(t)\le C+\displaystyle\int_{t_0}^t[f(s)x(s)+K]\dd s$，则 $x(t)\le[C+K(t-t_0)]\exp(\displaystyle\int_{t_0}^tf)$。}
\begin{proof}
设 $H=C+K(t-t_0)+\displaystyle\int_{t_0}^tf(s)x(s)\dd s$，$A(t)=\displaystyle\int_{t_0}^tf(s)\dd s$。由 $x\le H$ 得 $H'\le K+fH$。乘 $\e^{-A}$ 并积分：
\[
 H(t)\le\e^{A(t)}\left[C+K\int_{t_0}^t\e^{-A(s)}\dd s\right].
\]
因 $f\ge0$，$A(s)\ge0$，所以积分不超过 $t-t_0$。结合 $K\ge0$ 与 $x\le H$ 得所求估计。
这里 $H(t_0)=C$，所以积分因子步骤的初始项就是 $C$。最后的估计只把 $\displaystyle K\int_{t_0}^t\e^{-A(s)}\dd s$ 替为 $K(t-t_0)$，利用的是 $K\ge0$；若没有这个符号条件，不能保持相同方向地作该替换。
\end{proof}
\problem{5.5.1}{152}{设 $f,g$ 在平面区域 $G$ 连续并满足 Lipschitz 条件，$(x_0,y_0)\in G$。证明 $f(x_0,y_0)=g(x_0,y_0)=0$ 当且仅当系统 $x'=f(x,y),y'=g(x,y)$ 在该点任意邻域内都有时间长度任意大的轨道段。这里把解 $(x(t),y(t))$ 看作平面上以 $t$ 为参数的曲线，称为轨道。}
\begin{proof}
若两者均为零，常值解 $(x,y)=(x_0,y_0)$ 在任意长时间内都停留于每个邻域中，结论成立。

反之，若向量场在该点非零，至少一个分量非零。不妨 $f(x_0,y_0)>0$。由连续性，可取闭小矩形 $|x-x_0|\le r,|y-y_0|\le r$ 包含于 $G$，且其内 $f\ge c>0$。任何完全留在该矩形内的轨道段，在时间长度 $T$ 内都有 $x(t_2)-x(t_1)\ge cT$，而矩形内横坐标差至多 $2r$，所以 $T\le2r/c$。这与任意长轨道段矛盾。其他分量或负号情形只需反向比较相应坐标。
以 $f(x_0,y_0)>0$ 的情形说明，连续性允许取 $c=f(x_0,y_0)/2$ 并缩小矩形，使整个矩形内 $f\ge c$。若非零分量为负，改用 $-x$；若非零的是 $g$，改用 $y$ 或 $-y$。四种情况都得到同一有限停留时间上界，因而分类完备。
\end{proof}
\problem{5.5.3}{152}{对 $x'=\sin(tx)$，求解关于初始时刻及初值的导数在 $t_0=0,x_0=0$ 的表达式；若 $\varphi(t,\eta)$ 满足 $x(0)=\eta$，证明 $\partial\varphi/\partial\eta>0$。}
在 $t_0=0,x_0=0$ 时，零函数满足初值，且 $f_x=t\cos(tx)$ 连续，故由唯一性基准解就是 $x(t)=0$。初值导数 $U=\partial x/\partial x_0$ 满足变分方程
\[
U'=f_x(t,0)U=tU,\qquad U(0)=1.
\]
分离积分给 $\ln U=t^2/2$，所以 $U=\e^{t^2/2}$。

对初始时刻导数，使用 $\partial x/\partial t_0=-U(t)f(t_0,x_0)$。本题 $f(0,0)=\sin0=0$，故该导数为0。这里求的是初始时刻的偏导，不是 $\partial x/\partial t=f(t,x)$，两者的含义不同。
\ans{\left.\dfrac{\partial x}{\partial t_0}\right|_{t_0=x_0=0}=0,\qquad
\left.\dfrac{\partial x}{\partial x_0}\right|_{t_0=x_0=0}=\e^{t^2/2}.}
对一般 $\eta$，变分方程为 $U'=t\cos(t\varphi(t,\eta))U,U(0)=1$，所以
\[
 \dfrac{\partial\varphi}{\partial\eta}
 =\exp\left(\int_0^t s\cos(s\varphi(s,\eta))\dd s\right)>0.
\]
$f$ 有界使所有解可延拓到全实轴，上式对所有有限 $t$ 都有意义。
\problem{5.5.4}{152--153}{证明初值解 $x(t;t_0,x_0)$ 满足
\[
 \dfrac{\partial x}{\partial t_0}+\sum_{i=1}^n
 \dfrac{\partial x}{\partial x_0^i}f_i(t_0,x_0)=0.
\]}
\begin{proof}
采用本节可微依赖性定理的条件，使所写导数存在。由唯一性，先沿原解从 $t_0$ 走到 $s$，再从 $s$ 出发走到 $t$，得到同一终点：
\[
 x\bigl(t;s,x(s;t_0,x_0)\bigr)=x(t;t_0,x_0).
\]
在 $s=t_0$ 处求导。右端与 $s$ 无关，导数为零；左端由链式法则为
\[
 x_{t_0}(t;t_0,x_0)+D_{x_0}x(t;t_0,x_0)\,f(t_0,x_0).
\]
按矩阵乘法把第二项逐列展开，即为题述恒等式。这里 $x_{t_0}$ 与 $x_{x_0^i}$ 均为列向量，不能把分量下标与自变量混淆。
按分量展开链式法则，第 $j$ 个分量的导数为
\[
\dfrac{\partial x_j}{\partial t_0}
+\sum_{i=1}^n\dfrac{\partial x_j}{\partial x_0^i}
\left.\dfrac\dd{\dd s}x_i(s;t_0,x_0)\right|_{s=t_0}.
\]
最后一个导数由原微分方程等于 $f_i(t_0,x_0)$。对 $j=1,\ldots,n$ 集合这些标量式，便得到向量恒等式。
\end{proof}
\problem{5.6.2}{159}{编写 Euler、改进 Euler、四阶 Runge--Kutta 程序，以 $h=0.1$ 计算 $x'=x+\sin t,x(0)=1/2$，$0\le t\le1$，并与解析解 $x=\e^t-(\sin t+\cos t)/2$ 比较。}
先核对解析解。将方程乘 $\e^{-t}$，得到 $(x\e^{-t})'=\e^{-t}\sin t$，而
\[
\int\e^{-t}\sin t\dd t=-\dfrac12\e^{-t}(\sin t+\cos t).
\]
故 $x=C\e^t-(\sin t+\cos t)/2$，初值 $1/2=C-1/2$ 给 $C=1$。

现在三种算法都从 $t_0=0,X_0=0.5$、$h=0.1$ 开始。首步分别算如下。
\begin{enumerate}
\item Euler 法只有一个起点斜率 $f(0,0.5)=0.5$，所以 $X_1=0.5+0.1\cdot0.5=0.55$。
\item 改进 Euler 法先预测 $\widetilde X_1=0.55$，再计算预测终点的斜率 $f(0.1,0.55)=0.55+\sin0.1$。于是
\[
X_1=0.5+0.05[0.5+0.55+\sin0.1]\approx0.557491671.
\]
\item RK4 法的四个斜率逐个使用上一步的结果：
\[
\begin{aligned}
K_1&=0.5,\\
K_2&=0.5+0.05K_1+\sin0.05\approx0.574979169,\\
K_3&=0.5+0.05K_2+\sin0.05\approx0.578728128,\\
K_4&=0.5+0.1K_3+\sin0.1\approx0.657706229,\\
X_1&=0.5+\dfrac{0.1}{6}(K_1+2K_2+2K_3+K_4)
\approx0.557752014.
\end{aligned}
\]
\end{enumerate}
后续每步将 $t_n,X_n$ 代入相同的递推，必须分别保存三种算法各自的 $X_n$，不能混用。表中精确值逐点代入解析式，绝对误差用 $|X_n-x(t_n)|$ 计算；推进时保存完整精度，最后显示时再舍入。可运行程序见 \texttt{numerics.py}。
\begin{center}\small\setlength{\tabcolsep}{3pt}
\begin{tabular}{rrrrrrr}\toprule
$t$&精确值&Euler&绝对误差&改进Euler&绝对误差&RK4\\\midrule
0.0&0.500000000&0.500000000&$0$&0.500000000&$0$&0.500000000\\
0.1&0.557752127&0.550000000&$7.752\times10^{-3}$&0.557491671&$2.605\times10^{-4}$&0.557752014\\
0.2&0.632034804&0.614983342&$1.705\times10^{-2}$&0.631452601&$5.822\times10^{-4}$&0.632034550\\
0.3&0.724430460&0.696348609&$2.808\times10^{-2}$&0.723457947&$9.725\times10^{-4}$&0.724430036\\
0.4&0.836585029&0.795535490&$4.105\times10^{-2}$&0.835145560&$1.439\times10^{-3}$&0.836584403\\
0.5&0.970217220&0.914030874&$5.619\times10^{-2}$&0.968225130&$1.992\times10^{-3}$&0.970216352\\
0.6&1.127129756&1.053376515&$7.375\times10^{-2}$&1.124489297&$2.640\times10^{-3}$&1.127128603\\
0.7&1.309222770&1.215178414&$9.404\times10^{-2}$&1.305826893&$3.396\times10^{-3}$&1.309221285\\
0.8&1.518509528&1.401118024&$1.174\times10^{-1}$&1.514238495&$4.271\times10^{-3}$&1.518507656\\
0.9&1.757134672&1.612965435&$1.442\times10^{-1}$&1.751854467&$5.280\times10^{-3}$&1.757132352\\
1.0&2.027395183&1.852594670&$1.748\times10^{-1}$&2.020955715&$6.439\times10^{-3}$&2.027392347\\
\bottomrule\end{tabular}\end{center}
\begin{center}\small\begin{tabular}{rr}\toprule $t$&RK4绝对误差\\\midrule
0.1&$1.134\times10^{-7}$\\
0.2&$2.533\times10^{-7}$\\
0.3&$4.233\times10^{-7}$\\
0.4&$6.269\times10^{-7}$\\
0.5&$8.686\times10^{-7}$\\
0.6&$1.153\times10^{-6}$\\
0.7&$1.486\times10^{-6}$\\
0.8&$1.872\times10^{-6}$\\
0.9&$2.320\times10^{-6}$\\
1.0&$2.836\times10^{-6}$\\
\bottomrule\end{tabular}\end{center}
\ans{X_{10}^{\mathrm E}=1.852594670,\quad X_{10}^{\mathrm H}=2.020955715,\quad X_{10}^{\mathrm {RK4}}=2.027392347.}
\problem{5.6.3}{159}{取 $h=0.1$，用四阶 Runge--Kutta 法计算 $\displaystyle\int_0^1\e^{-t^2}\dd t$。}
把积分转为初值问题 $u'=f(t)=\e^{-t^2},u(0)=0$，这样 $u(1)$ 就是目标积分。因右端不依赖 $u$，每一步 RK4 中
\[
K_1=f(t_n),\quad K_2=K_3=f(t_n+h/2),\quad K_4=f(t_n+h).
\]
所以其加权平均恰为
\[
u_{n+1}-u_n=\dfrac h6[f(t_n)+4f(t_n+h/2)+f(t_n+h)].
\]
首步 $t_0=0,h=0.1$，三个不同取样值为
\[
1,\quad\e^{-0.0025}\approx0.997503122397,\quad
\e^{-0.01}\approx0.990049833749.
\]
加权后首步增量为 $[1+4\e^{-0.0025}+\e^{-0.01}]/60\approx0.099667671389$。将起点依次换为 $0,0.1,\ldots,0.9$，得到下面10个增量；从 $u_0=0$ 累加所有增量就是最终近似值。表中每行只展示舍入结果，最终和使用未舍入的数值计算。
\begin{center}\small\begin{tabular}{rr}\toprule 区间& RK4积分增量\\\midrule
$[0.0,0.1]$&0.099667705389\\
$[0.1,0.2]$&0.097697403695\\
$[0.2,0.3]$&0.093872881261\\
$[0.3,0.4]$&0.088414976567\\
$[0.4,0.5]$&0.081628175041\\
$[0.5,0.6]$&0.073872517703\\
$[0.6,0.7]$&0.065532128959\\
$[0.7,0.8]$&0.056984168619\\
$[0.8,0.9]$&0.048571634515\\
$[0.9,1.0]$&0.040582592127\\
\bottomrule\end{tabular}\end{center}
\ans{I_{\mathrm{RK4}}=0.7468241838759147,\qquad I_{\mathrm{ref}}=0.746824132812427,\\ |I_{\mathrm{RK4}}-I_{\mathrm{ref}}|\approx 5.106\times10^{-8}.}
\checknote{该递推与复合 Simpson 求积在20个半步小区间上的权重一致。另用高精度积分独立计算参考值，误差见结果；这里的 $h=0.1$ 是 RK4 的整步长。}
